\scanpage{22}
% L1–2
\uline{Corollaire 3} On a un hom naturel de faisceaux
\[
\underline S(\overline X/X_0)\longrightarrow\underline{\operatorname{Is}}_{X_0}(\overline{\deleted{X}\xi},\overline X)
\]
% L3–7
compatible avec les opérations de $\underline{\operatorname{Aut}}_T^{X_0}(X)$ à gauche, et les opérations \unsure{---|à d.} de $\underline{\operatorname{Aut}}_T^{X_0}(\xi)$ [qui opère sur le premier \unsure{mb|m$^{\rm bre}$} par le corollaire 2]. Cet homomorphisme est \uline{injectif} [car $\underline{\operatorname{Aut}}_T^{X_0}(\xi)$ opère fidèlement dans $\overline\xi$]

% L8–12
On peut dire aussi [mais cela n'est pas très \unsure{malin|utile}] que $\overline X/X_0$ \uline{est muni d'une structure de} \deleted{fais} \uline{schéma sur} $X_0$ \uline{à faisceau structural} $\unsure{=|:}\boxed{\underline{\operatorname{Aut}}_T^{X_0}(\xi)=\underline G_{X_0}^{\underline{\Lambda}}}$ \uline{opérant sur le schéma-type} $\xi=X_0\times_S T=X_0\otimes_{\underline{O}_S}\underline{\Lambda}$

% L13–16
\uline{La donnée de} $X$ \uline{simple sur} $T$, se \uline{réduisant suivant} $X_0$, est \uuline{équivalente} \uline{à la donnée d'un schéma} $\overline X$ \uline{sur} $X_0$, \uline{muni d'un tel faisceau structural}.

% L17–18
\uline{Le faisceau principal associé n'est autre que} $\underline S(\overline X/X_0)$ !

% L19–21
$\overline\xi$ est noté aussi $X_0^{\{\Lambda\}}$, et appelé \deleted{\unsure{S-schéma}} \uline{$S$-préschéma des p$^{\rm ts}$ proches} de $X_0$ \uline{du type} $\Lambda$. \unsure{Ses|les} sections \deleted{\unsure{sont des}} forment donc
% archive sheet number: 21

